% GENERATED FILE. Edit canonical lessons, then run npm run generate:books.
% canonical-source-hash: fnv1a64:e6632ed7a4ebfd06
% canonical-lessons: TA-C389-ceyalpatuttu, TA-C389-nataimurai, TA-C389-valimurai, TA-C389-natavatikkai, TA-C389-nirvaki

\chapter{From Plan to Practice}
\label{ch:ta-from-plan-to-practice}
\hlchaptermodality{389}

\begin{chapteropening}
\textbf{By the end of this chapter:} \emph{I can say that a plan is put into practice and managed.}

Complete the last lesson of chapter 389: I can say that a plan is put into practice and managed.
\end{chapteropening}

\section[\texorpdfstring{\ta{செயல்படுத்து} (ceyalpaṭuttu)}{செயல்படுத்து (ceyalpaṭuttu)}]{\ta{செயல்படுத்து} (ceyalpaṭuttu) — to implement; to put into action}

\label{lesson:TA-C389-ceyalpatuttu}



\begin{quote}
Before the new one: say the Tamil for to inform, then the Tamil for a session.
\end{quote}

\subsection*{You'll want to know: \ta{செயல்படுத்து}}
\textbf{\ta{செயல்படுத்து}} — \emph{ceyalpaṭuttu} — \textquotedblleft{}to implement; to put into action\textquotedblright{}.

\begin{quote}
\ta{அரசு} \ta{இந்தத்} \ta{திட்டத்தை} \ta{அடுத்த} \ta{ஆண்டு} \ta{செயல்படுத்தும்}.
\end{quote}

\begin{quote}
\emph{aracu intat tiṭṭattai aṭutta āṇṭu ceyalpaṭuttum.}
\end{quote}

\begin{quote}
The government will implement this plan next year.
\end{quote}

Past (I …): \ta{செயல்படுத்தினேன்}, \emph{ceyalpaṭuttiṉēṉ}.

\ta{செயல்}, action + \ta{படுத்து}, to make: to make a plan act.

\begin{cousinweb}
\textbf{\ta{செயல்படுத்து}} (\emph{ceyalpaṭuttu}) is native: \emph{ceyal}, \textquotedblleft{}act, action, work\textquotedblright{}, from \emph{cey}, \textquotedblleft{}to do\textquotedblright{} (DEDR 1957), with \emph{paṭuttu}, the causative of \emph{paṭu}.
\end{cousinweb}

\begin{grammarlens}[title={What you've built}]
One more word for saying what happens if.
\end{grammarlens}

\subsection*{Your turn}
\begin{itemize}
\raggedright
  \item \emph{Say it:} \emph{ceyalpaṭuttu}
  \item \emph{Say it:} \emph{ceyalpaṭuttu}, once more
  \item \emph{Say it:} \emph{aracu intat tiṭṭattai aṭutta āṇṭu ceyalpaṭuttum.} — The government will implement this plan next year.
  \item \emph{Recall:} say the Tamil for to inform, then the Tamil for a session, then say \emph{ceyalpaṭuttu} again
\end{itemize}

\begin{tcolorbox}[breakable,colback=teal!4,colframe=teal!35!black,title={Before you move on}]
What does \ta{செயல்படுத்து} mean? (\textquotedblleft{}To implement; to put into action\textquotedblright{}.) Say it once more.
\end{tcolorbox}
\section[\texorpdfstring{\ta{நடைமுறை} (naṭaimuṟai)}{நடைமுறை (naṭaimuṟai)}]{\ta{நடைமுறை} (naṭaimuṟai) — practice; procedure}

\label{lesson:TA-C389-nataimurai}



\begin{quote}
Before the new one: say the Tamil for a session, then the Tamil for to implement.
\end{quote}

\subsection*{You'll want to know: \ta{நடைமுறை}}
\textbf{\ta{நடைமுறை}} — \emph{naṭaimuṟai} — \textquotedblleft{}practice; procedure\textquotedblright{}.

\begin{quote}
\ta{புதிய} \ta{விதி} \ta{நாளை} \ta{முதல்} \ta{நடைமுறைக்கு} \ta{வரும்}.
\end{quote}

\begin{quote}
\emph{putiya viti nāḷai mutal naṭaimuṟaikku varum.}
\end{quote}

\begin{quote}
The new rule comes into force from tomorrow.
\end{quote}

\ta{நடைமுறைக்கு} \ta{வா}, to come into force; \ta{நடைமுறையில்}, in practice.

\begin{cousinweb}
\textbf{\ta{நடைமுறை}} (\emph{naṭaimuṟai}) is native: \emph{naṭai}, \textquotedblleft{}walk, conduct\textquotedblright{}, from \emph{naṭa} (DEDR 3582), and \emph{muṟai}, \textquotedblleft{}order, arrangement, way\textquotedblright{} (DEDR 5015): the way things walk.
\end{cousinweb}

\begin{grammarlens}[title={What you've built}]
One more word for saying what happens if.
\end{grammarlens}

\subsection*{Your turn}
\begin{itemize}
\raggedright
  \item \emph{Say it:} \emph{naṭaimuṟai}
  \item \emph{Say it:} \emph{naṭaimuṟai}, once more
  \item \emph{Say it:} \emph{putiya viti nāḷai mutal naṭaimuṟaikku varum.} — The new rule comes into force from tomorrow.
  \item \emph{Recall:} say the Tamil for a session, then the Tamil for to implement, then say \emph{naṭaimuṟai} again
\end{itemize}

\begin{tcolorbox}[breakable,colback=teal!4,colframe=teal!35!black,title={Before you move on}]
What does \ta{நடைமுறை} mean? (\textquotedblleft{}Practice; procedure\textquotedblright{}.) Say it once more.
\end{tcolorbox}
\section[\texorpdfstring{\ta{வழிமுறை} (vaḻimuṟai)}{வழிமுறை (vaḻimuṟai)}]{\ta{வழிமுறை} (vaḻimuṟai) — a procedure; a method}

\label{lesson:TA-C389-valimurai}



\begin{quote}
Before the new one: say the Tamil for to implement, then the Tamil for practice.
\end{quote}

\subsection*{You'll want to know: \ta{வழிமுறை}}
\textbf{\ta{வழிமுறை}} — \emph{vaḻimuṟai} — \textquotedblleft{}a procedure; a method\textquotedblright{}.

\begin{quote}
\ta{விண்ணப்பிக்கும்} \ta{வழிமுறை} \ta{மிக} \ta{எளிது}.
\end{quote}

\begin{quote}
\emph{viṇṇappikkum vaḻimuṟai mika eḷitu.}
\end{quote}

\begin{quote}
The procedure for applying is very simple.
\end{quote}

Plural: \ta{வழிமுறைகள்}, \emph{vaḻimuṟaikaḷ}.

Plural \ta{வழிமுறைகள்}, guidelines.

\begin{cousinweb}
\textbf{\ta{வழிமுறை}} (\emph{vaḻimuṟai}) is native: \emph{vaḻi}, \textquotedblleft{}way, means, method\textquotedblright{} (DEDR 5297), and \emph{muṟai}, \textquotedblleft{}order, system\textquotedblright{} (DEDR 5015).
\end{cousinweb}

\begin{grammarlens}[title={What you've built}]
One more word for saying what happens if.
\end{grammarlens}

\subsection*{Your turn}
\begin{itemize}
\raggedright
  \item \emph{Say it:} \emph{vaḻimuṟai}
  \item \emph{Say it:} \emph{vaḻimuṟai}, once more
  \item \emph{Say it:} \emph{viṇṇappikkum vaḻimuṟai mika eḷitu.} — The procedure for applying is very simple.
  \item \emph{Recall:} say the Tamil for to implement, then the Tamil for practice, then say \emph{vaḻimuṟai} again
\end{itemize}

\begin{tcolorbox}[breakable,colback=teal!4,colframe=teal!35!black,title={Before you move on}]
What does \ta{வழிமுறை} mean? (\textquotedblleft{}A procedure; a method\textquotedblright{}.) Say it once more.
\end{tcolorbox}
\section[\texorpdfstring{\ta{நடவடிக்கை} (naṭavaṭikkai)}{நடவடிக்கை (naṭavaṭikkai)}]{\ta{நடவடிக்கை} (naṭavaṭikkai) — a course of action; a measure}

\label{lesson:TA-C389-natavatikkai}



\begin{quote}
Before the new one: say the Tamil for practice, then the Tamil for a procedure.
\end{quote}

\subsection*{You'll want to know: \ta{நடவடிக்கை}}
\textbf{\ta{நடவடிக்கை}} — \emph{naṭavaṭikkai} — \textquotedblleft{}a course of action; a measure\textquotedblright{}.

\begin{quote}
\ta{இந்தப்} \ta{பிரச்சினைக்கு} \ta{அரசு} \ta{உடனே} \ta{நடவடிக்கை} \ta{எடுக்க} \ta{வேண்டும்}.
\end{quote}

\begin{quote}
\emph{intap piracciṉaikku aracu uṭaṉē naṭavaṭikkai eṭukka vēṇṭum.}
\end{quote}

\begin{quote}
The government must take action on this problem at once.
\end{quote}

\ta{நடவடிக்கை} \ta{எடு}, to take action, to take steps.

\begin{cousinweb}
\textbf{\ta{நடவடிக்கை}} (\emph{naṭavaṭikkai}) is native Dravidian: DEDR 3582 lists it in the family of \emph{naṭa}, \textquotedblleft{}to walk, proceed, happen\textquotedblright{}, meaning \textquotedblleft{}conduct, proceeding\textquotedblright{}.
\end{cousinweb}

\begin{grammarlens}[title={What you've built}]
One more word for saying what happens if.
\end{grammarlens}

\subsection*{Your turn}
\begin{itemize}
\raggedright
  \item \emph{Say it:} \emph{naṭavaṭikkai}
  \item \emph{Say it:} \emph{naṭavaṭikkai}, once more
  \item \emph{Say it:} \emph{intap piracciṉaikku aracu uṭaṉē naṭavaṭikkai eṭukka vēṇṭum.} — The government must take action on this problem at once.
  \item \emph{Recall:} say the Tamil for practice, then the Tamil for a procedure, then say \emph{naṭavaṭikkai} again
\end{itemize}

\begin{tcolorbox}[breakable,colback=teal!4,colframe=teal!35!black,title={Before you move on}]
What does \ta{நடவடிக்கை} mean? (\textquotedblleft{}A course of action; a measure\textquotedblright{}.) Say it once more.
\end{tcolorbox}
\section[\texorpdfstring{\ta{நிர்வகி} (nirvaki)}{நிர்வகி (nirvaki)}]{\ta{நிர்வகி} (nirvaki) — to administer; to manage}

\label{lesson:TA-C389-nirvaki}



\begin{quote}
Before the new one: say the Tamil for a procedure, then the Tamil for a course of action.
\end{quote}

\subsection*{You'll want to know: \ta{நிர்வகி}}
\textbf{\ta{நிர்வகி}} — \emph{nirvaki} — \textquotedblleft{}to administer; to manage\textquotedblright{}.

\begin{quote}
\ta{என்} \ta{அம்மா} \ta{எங்கள்} \ta{குடும்பக்} \ta{கடையை} \ta{நிர்வகிக்கிறார்}.
\end{quote}

\begin{quote}
\emph{eṉ ammā eṅkaḷ kuṭumpak kaṭaiyai nirvakikkiṟār.}
\end{quote}

\begin{quote}
My mother manages our family shop.
\end{quote}

Past (I …): \ta{நிர்வகித்தேன்}, \emph{nirvakittēṉ}.

The noun is \ta{நிர்வாகம்}, administration; \ta{நிர்வாகி}, an administrator.

\begin{cousinweb}
\textbf{\ta{நிர்வகி}} (\emph{nirvaki}) is a verb built on Sanskrit \emph{nirvāha}, \textquotedblleft{}carrying on, accomplishing, completion\textquotedblright{} (Monier-Williams). Tamil writes the h between vowels as \ta{க}.
\end{cousinweb}

\begin{grammarlens}[title={What you've built}]
One more word for saying what happens if.
\end{grammarlens}

\subsection*{Your turn}
\begin{itemize}
\raggedright
  \item \emph{Say it:} \emph{nirvaki}
  \item \emph{Say it:} \emph{nirvaki}, once more
  \item \emph{Say it:} \emph{eṉ ammā eṅkaḷ kuṭumpak kaṭaiyai nirvakikkiṟār.} — My mother manages our family shop.
  \item \emph{Recall:} say the Tamil for a procedure, then the Tamil for a course of action, then say \emph{nirvaki} again
\end{itemize}

\begin{tcolorbox}[breakable,colback=teal!4,colframe=teal!35!black,title={Before you move on}]
What does \ta{நிர்வகி} mean? (\textquotedblleft{}To administer; to manage\textquotedblright{}.) Say it once more.
\end{tcolorbox}
